% ---------------------------------------------------------------------------
% CITC thesis / capstone template.
% Build:  latexmk -pdf thesis.tex        (biber runs automatically)
% ---------------------------------------------------------------------------
\documentclass[12pt,a4paper]{article}

\usepackage{graphicx}
\usepackage{amsmath}
\usepackage{amssymb}
\usepackage{booktabs}
\usepackage{float}
\usepackage[colorlinks=true,allcolors=black]{hyperref}
\usepackage{citc}          % all CITC formatting; load after hyperref
\usepackage{cleveref}      % optional; citc.sty teaches it "Chapter"/"Section"

\title{Your Thesis Title}
\author{Author One \and Author Two}
\date{}

\begin{document}
\maketitle

% ===========================================================================
% A \section is a CHAPTER: it starts a new page and prints as
%     Chapter I
%     INTRODUCTION
% centered, bold, uppercase. Use the preferred titles from Section 4.3.1:
% INTRODUCTION, REVIEW OF LITERATURE, METHODOLOGY, RESULTS AND DISCUSSION,
% SUMMARY AND CONCLUSION, RECOMMENDATIONS.
% ===========================================================================
\section{Introduction}
\label{sec:intro}

\subsection{Background of the Study}          % major subsection  -> 1.1.
\label{sec:background}

Each paragraph is indented 1.25\,cm, justified and single-spaced. The first
paragraph after a heading is indented too.

\subsubsection{A Minor Subsection}            % minor subsection  -> 1.1.1.

\phead{A paragraph heading} It is indented, bold and italicized, in sentence
case, followed by a period, and the paragraph runs on from it.

\subsection{Statement of the Problem}
\subsection{Objectives of the Study}
\subsection{Scope and Limitations of the Study}

% ===========================================================================
\section{Review of Literature}
\label{sec:rrl}

Cite in APA~7 through \textsf{biblatex}: parenthetically with
\verb|\parencite| \parencite{daganzo1994}, or narratively with
\verb|\textcite|, as \textcite{daganzo1994} does. In a table column that
already names the author, use \verb|\workcite{Name}{key}| so the name is not
printed twice.

% ---------------------------------------------------------- table (Article 5)
% Title ABOVE, sentence case, left aligned, label "Table X.X.".
% First-level headings UPPERCASE. Single rules top and bottom, no vertical
% rules, no side boxes. tabularx keeps a wide column inside the margins.
\begin{table}[H]
\centering
\small
\caption{A short sentence-case title that names the table.}
\label{tab:example}
\begin{tabularx}{\linewidth}{@{}p{3.2cm}X@{}}
\toprule
FIRST LEVEL HEADING & FIRST LEVEL HEADING \\
\midrule
Row heading 1 & Entry \\
Row heading 2 & Entry \\
\bottomrule
\end{tabularx}
\end{table}

% --------------------------------------------------------- figure (Article 6)
% Caption BELOW, sentence case, left aligned, label "Figure Y.Y.".
% Keep the caption a NAME; put the description in the text.
\begin{figure}[htbp]
\centering
\rule{6cm}{3cm}                      % replace with \includegraphics
\caption{A short name for the figure.}
\label{fig:example}
\end{figure}

% ===========================================================================
\section{Methodology}
\label{sec:method}

% ------------------------------------------------------- equation (Article 7)
% Numbered <chapter>-<n>, flush right as "(Equation 3-1)", followed by a
% "where:" list naming every term, with units in parentheses.
\begin{equation}
y_i = \min\bigl(D_i,\; S_i,\; C_i\bigr)
\label{eq:example}
\end{equation}
\begin{whereterms}
$y_i$ & is the flow across boundary $i$ (veh/h) \\
$D_i$ & is the demand (veh/h) \\
$S_i$ & is the supply (veh/h) \\
$C_i$ & is the capacity (veh/h)
\end{whereterms}

Refer to numbered material as \Cref{eq:example}, \Cref{tab:example},
\Cref{fig:example}, \Cref{sec:rrl}.

\printbibliography[title={REFERENCES}]

\end{document}
